\begin{otherlanguage}{italian}
\chapter*{Sommario}
\addcontentsline{toc}{chapter}{Sommario}

\begin{noparskip}
Per comprendere il motivo per cui le reti neurali addestrate con il metodo del gradiente ottengono risultati positivi, sono necessari modelli teorici in grado di cogliere il modo in cui esse individuano strutture a bassa dimensionalità nascoste all'interno di dati ad alta dimensionalità. Nonostante il predominio empirico del deep learning su larga scala, anche le reti di piccole dimensioni e poco profonde rimangono analiticamente sfuggenti: le dinamiche che regolano ciò che apprendono, e la velocità con cui lo fanno, non sono ancora pienamente comprese.
La presente tesi si basa su una teoria di tali dinamiche ispirata alla fisica statistica: studia la discesa stocastica del gradiente su reti poco profonde e su architetture di attenzione semplificate, addestrate su dati sintetici generati da un «insegnante», monitorando l'evoluzione di un ristretto numero di statistiche riassuntive anziché di ogni singolo peso.
Nel limite ad alta dimensionalità, tale evoluzione è descritta esattamente da un sistema chiuso di equazioni differenziali deterministiche, rendendo la velocità e la difficoltà dell'apprendimento pienamente quantificabili.

Inizialmente, inquadriamo descrizioni di questa dinamica precedentemente separate — i regimi classici del flusso di gradiente, ad alta dimensionalità e del campo medio sovraparametrizzato — come casi particolari di un unico processo sottostante. All'interno di questo quadro, forniamo formule in forma chiusa relative alla complessità campionaria per l'apprendimento di obiettivi difficili, del tipo «recupero di fase»~(phase retrieval), e dimostriamo che l'aumento della larghezza della rete accelera l'apprendimento solo di un fattore costante, senza sfuggire alla difficoltà determinata dall'esponente di informazione dell'obiettivo.

La seconda parte studia come semplici modifiche alla discesa stocastica del gradiente standard (vanilla stochastic gradient descent) consentano di spostare l'addestramento al di fuori di questo panorama di difficoltà. Un modesto riutilizzo dei dati, ovvero l'osservazione di ciascun campione due volte anziché una sola, sposta l'ostacolo dall'\emph{esponente di informazione}~(information exponent) che governa l'apprendimento a passaggio singolo verso un \emph{esponente generativo}~(generative exponent) più piccolo, consentendo di recuperare in modo efficiente quasi tutte le direzioni nascoste, con un meccanismo gerarchico che tiene conto dei casi difficili rimanenti. Identifichiamo la dimensione del batch che minimizza il tempo di addestramento per un obiettivo di data difficoltà e introduciamo una funzione di perdita modificata che elimina del tutto il limite di velocità risultante.

Infine, estendiamo l'analisi oltre gli input statici ai dati sequenziali e a un livello di attenzione semplificato, dimostrando che la struttura posizionale può accelerare sostanzialmente l'apprendimento rispetto agli input non strutturati.

Nel loro insieme, questi risultati forniscono una teoria risolvibile e unificata delle dinamiche di addestramento delle reti poco profonde e dei modelli di attenzione semplici. Essi quantificano quando l'apprendimento basato sui gradienti riesce a scoprire la struttura nascosta, fino a che punto semplici modifiche algoritmiche possano potenziare tale successo e quali proprietà dei dati e dell'addestramento ne determinino la velocità.

\end{noparskip}
\textbf{Parole chiave:} apprendimento profondo, reti neurali, discesa del gradiente stocastica, dinamiche di addestramento, fisica statistica, teoria di campo medio, modelli insegnante-studente (teacher-student), complessità campionaria, recupero di fase, meccanismi di attenzione
\end{otherlanguage}
